\documentclass{beamer}
\usetheme{Madrid} % You can change this theme if you like
\usepackage{amsmath} % For math equations

\title[Labor Market Exercises]{Numerical Problems and Exercises: Labor Market and Unemployment}
\author{Economics Faculty}
\date{\today}

\begin{document}

% Title Slide
\begin{frame}
    \titlepage
\end{frame}

% Table of Contents
\begin{frame}{Overview}
    \tableofcontents
\end{frame}

% Section 1: Numerical Problems
\section{Numerical Problems}

% Problem 1: Calculating the Unemployment Rate
\subsection{Problem 1: Calculating the Unemployment Rate}
\begin{frame}{Problem 1: Calculating the Unemployment Rate}
The labor force in a country consists of 150 million people, of which 10 million are unemployed.
\begin{enumerate}
    \item Calculate the unemployment rate.  
    \item If 2 million unemployed people stop looking for jobs, what happens to the unemployment rate?  
\end{enumerate}
\end{frame}

\begin{frame}{Answer to Problem 1}
(a) Unemployment Rate:
\[
\text{Unemployment Rate} = \frac{\text{Unemployed}}{\text{Labor Force}} \times 100 = \frac{10}{150} \times 100 = 6.67\%
\]

(b) Adjusted Unemployment Rate:  
If 2 million unemployed stop looking for jobs, the labor force becomes \(150 - 2 = 148 \, \text{million}\), and the unemployment rate becomes:
\[
\text{Unemployment Rate} = \frac{8}{148} \times 100 \approx 5.41\%
\]
\end{frame}

% Problem 2: Labor Force Participation Rate
\subsection{Problem 2: Labor Force Participation Rate}
\begin{frame}{Problem 2: Labor Force Participation Rate}
The working-age population of a country is 200 million. Of these, 140 million are in the labor force, and 60 million are not participating in the labor market.
\begin{enumerate}
    \item Calculate the labor force participation rate.  
\end{enumerate}
\end{frame}

\begin{frame}{Answer to Problem 2}
\[
\text{Labor Force Participation Rate} = \frac{\text{Labor Force}}{\text{Working-Age Population}} \times 100
\]
Substituting values:
\[
\text{Labor Force Participation Rate} = \frac{140}{200} \times 100 = 70\%
\]
\end{frame}

% Problem 3: Employment Rate
\subsection{Problem 3: Employment Rate}
\begin{frame}{Problem 3: Employment Rate}
In a small country, the labor force is 25 million, and 22 million people are employed.
\begin{enumerate}
    \item Calculate the unemployment rate.  
    \item Calculate the employment rate.
\end{enumerate}
\end{frame}

\begin{frame}{Answer to Problem 3}
(a) Unemployment Rate:
\[
\text{Unemployed} = 25 - 22 = 3 \, \text{million}
\]
\[
\text{Unemployment Rate} = \frac{3}{25} \times 100 = 12\%
\]

(b) Employment Rate:
\[
\text{Employment Rate} = \frac{\text{Employed}}{\text{Labor Force}} \times 100 = \frac{22}{25} \times 100 = 88\%
\]
\end{frame}

% Problem 4: Cyclical Unemployment
\subsection{Problem 4: Cyclical Unemployment}
\begin{frame}{Problem 4: Cyclical Unemployment}
The natural rate of unemployment in a country is estimated to be 4\%. During a recession, the actual unemployment rate rises to 8\%.
\begin{enumerate}
    \item What is the cyclical unemployment rate?  
    \item If the labor force is 50 million, how many people are cyclically unemployed?  
\end{enumerate}
\end{frame}

\begin{frame}{Answer to Problem 4}
(a) Cyclical Unemployment Rate:
\[
\text{Cyclical Unemployment Rate} = \text{Actual Unemployment Rate} - \text{Natural Unemployment Rate}
\]
\[
\text{Cyclical Unemployment Rate} = 8\% - 4\% = 4\%
\]

(b) Number of Cyclically Unemployed:
\[
\text{Cyclically Unemployed} = 4\% \times 50 \, \text{million} = 0.04 \times 50 = 2 \, \text{million}
\]
\end{frame}

% Problem 5: Structural vs. Frictional Unemployment
\subsection{Problem 5: Structural vs. Frictional Unemployment}
\begin{frame}{Problem 5: Structural vs. Frictional Unemployment}
In an economy, the total number of unemployed people is 8 million. Of these, 3 million are frictionally unemployed, and 2 million are structurally unemployed.
\begin{enumerate}
    \item What percentage of total unemployment is structural?  
    \item What percentage is frictional?  
    \item What portion of unemployment is due to other factors (e.g., cyclical)?  
\end{enumerate}
\end{frame}

\begin{frame}{Answer to Problem 5}
(a) Structural Unemployment Percentage:
\[
\text{Structural Unemployment} = \frac{2}{8} \times 100 = 25\%
\]

(b) Frictional Unemployment Percentage:
\[
\text{Frictional Unemployment} = \frac{3}{8} \times 100 = 37.5\%
\]

(c) Other Unemployment:
\[
\text{Other Unemployment} = 8 - (3 + 2) = 3 \, \text{million}
\]
\[
\text{Percentage} = \frac{3}{8} \times 100 = 37.5\%
\]
\end{frame}

% Discussion Questions
\section{Discussion Questions}
\begin{frame}{Discussion Questions}
\begin{enumerate}
    \item In an economy with high structural unemployment, should policymakers focus on demand-side or supply-side policies? Justify your answer.
    \item How does the inclusion of discouraged workers impact the accuracy of the unemployment rate?
    \item Compare the impact of automation on frictional and structural unemployment. Which one is likely to be affected more, and why?
\end{enumerate}
\end{frame}

% Closing Slide
\begin{frame}
    \centering
    \Huge{Thank You!}
\end{frame}

\end{document}

